\section{Hunting Iguanas}

Justice is the condition in which everything is in its proper place. Justice is the reflection of wisdom, which is necessary to understand the proper and right place for a thing or a being to be. For example, in nature, each animal species is in its ecological niche. It has adapted to that niche; it has food sources and natural enemies.

When a species is introduced into an environment which is not natural to it, it either cannot survive in it or else it creates an imbalance in the natural order of things. Reptilian pets like pythons and iguanas have created a problem in Florida so that iguanas are among invasive species now banned in Florida\footnote{\url{https://www.theguardian.com/us-news/2021/mar/22/toilet-invading-iguanas-invasive-species-banned-florida}}.

Pet owners let them out into the wild when they grow too large to keep at home. Pythons thrive in the Everglades but iguanas, like rats, do well in human environments. They are messy and replace native species of lizards; on the rare days when it gets cold, they fall out of trees. Although some people try to rescue them, iguanas are legal, and even encouraged to hunt down. The humane method is to use an air gun.

Iguanas are marvels of evolution: they are nimble afoot, and they can swim, jump, and climb. That makes them hard to shoot. We tried traps with enticing canned iguana food. One iguana was so stupid he couldn't find the trap door to get at the food. Instead, he kept bumping into the cage; discouraged he scampered away.

\paragraph{Right Action}


If justice is the condition of being in the proper place, then its complement, right action, is the cognitive action that actualizes the condition of justice. So if iguanas are not in their proper place, then the right action is to hunt them down.

In \emph{Prolegomena to the Metaphysics of Islam}, \textbf{Syed Naquib} calls right action \emph{adab}. It is a marvelous concept, since it relates seemingly disparate situations under the same theme of the condition of justice. He lists the following examples of adab:

\begin{itemize}


\item \textbf{Adab toward one's self}: The human self or soul has two aspects: the one predisposed to praiseworthy acts, intelligent by nature, loyal to its covenant with God; the other inclined to evil deeds, bestial by nature, heedless of its covenant with God. The former we call the rational soul the latter the carnal or animal soul, When the rational soul subdues the animal soul and renders it under control then one has put the animal soul in its proper place and the rational soul also in its proper place. In this way, and in relation to one's self, one is putting one's self in one's proper place.

\item \textbf{Adab toward family}: In relation to one's family and its various members, one's attitude and behaviour toward one's parents and elders display sincere acts of humility, love, respect, care, charity. This shows one knows one's proper place in relation to them by putting them in their proper places.

\item \textbf{Adab toward social circle}: Similarly, such attitude and behaviour, when extended to teachers, friends, community, leaders, manifest knowledge of one's proper place in relation to them; and this knowledge entails requisite acts in order to actualize adab toward them all.

\item \textbf{Adab toward language}: This is putting words in their proper places so that their true meanings become intelligible. Sentences and verses in like manner so that prose and poetry become literature.

\item \textbf{Adab toward nature and the environment}: When one puts trees and stones, mountains, rivers, valleys and lakes, animals and their habitat in their proper places.

\item \textbf{Adab toward home and furniture}: The same applies to one's home when one arranges furniture and puts things in their proper places therein until harmony is achieved.

\item \textbf{Adab toward art and music}: Putting colours, shapes, and sounds in their proper places producing pleasing effects.

\item \textbf{Adab toward knowledge}: Knowledge too, and its many branches and disciplines, some of which have more important bearing upon our life and destiny than others, grading them according to various levels and priorities and classifying the various sciences in relation to their priorities putting each one of them in its proper place.
\end{itemize}


Self-mastery, literature, art, and even keeping a well-designed home all fall under the same rubric.

\paragraph{Corruption of Right Order}


Right order takes the form of a hierarchy. Corruption of right order is marked when no one is in his right place. Places are determined by power, wealth, lineage, or reputation rather than the highest criteria of knowledge, intelligence, and virtue. False leaders emerge in all spheres of life. This comes about through the loss of knowledge. Even worse, there is not even the capacity and ability to even recognize and acknowledge true leaders. Naquib goes further:

\begin{quotex}
Because of the intellectual anarchy that characterizes this situation, the common people become determiners of intellectual decisions and are raised to the level of authority on matters of knowledge. Authentic definitions become undone, and in their stead we are left with platitudes and vague slogans disguised as profound concepts. The inability to define; to identify and isolate problems, and hence to provide for right solutions; the creation of pseudo-problems; the reduction of problems to mere political, socio-economic and legal factors become evident. It is not surprising if such a situation provides a fertile breeding ground for the emergence of deviationists and extremists of many kinds who make ignorance their capital.
\end{quotex}



\flrightit{Posted on 2021-09-26 by Cologero}

\begin{center}* * *\end{center}

\begin{footnotesize}
\begin{sffamily}
\texttt{scyld on 2021-09-29 at 21:53 said:}

how does rational soul conquer the animal soul?

\hfill

\texttt{Cologero on 2021-09-29 at 22:52 said:}

How does one trap and kill an iguana?

\hfill

\texttt{.chris on 2022-09-26 at 18:49 said:}

In being, brazenly, said sad Iguana, ignorant of its relevance.


\hfill

\end{sffamily}
\end{footnotesize}

